\documentclass[11pt]{article}
\usepackage{amsmath,amssymb,amsfonts}
\usepackage{geometry}
\usepackage{enumitem}
\usepackage{color}
\usepackage{hyperref}
\usepackage{booktabs}

\geometry{letterpaper, margin=1in}

\definecolor{addressed}{rgb}{0,0.5,0}
\definecolor{partial}{rgb}{0.8,0.4,0}
\definecolor{pending}{rgb}{0.8,0,0}

\title{\textbf{Response to Reviewers}\\[0.5em]
\large \textbf{New Title}: Flux-Preserving Adaptive Finite State Projection\\
for Multiscale Stochastic Reaction Networks\\[0.5em]
\textbf{Old Title: } Adaptive Finite State Projection with Quantile-Based\\
 Pruning for Solving the Chemical Master Equation\\[0.5em]
\normalsize Submitted to: SIAM Multiscale Modeling and Simulation}

\author{Aditya Dendukuri, Shivkumar Chandrasekaran, and Linda Petzold\\
Department of Computer Science\\
University of California, Santa Barbara}

\date{\today}

\begin{document}

\maketitle

\noindent Dear Editor and Reviewers,

\medskip

Thank you for the careful and very helpful evaluation of our manuscript. We have
undertaken a substantial revision in response to your comments, and both the scope
and the technical core of the paper have changed accordingly.

The original submission centered on a quantile-based pruning strategy: at each
time step, a fixed fraction of probability mass was removed and the distribution
was renormalized. As Reviewer~1 rightly pointed out, pruning based on probability
alone can mishandle bottleneck states, and the earlier error analysis was not
fully rigorous. In the revised manuscript we now present a flux-based adaptive FSP
framework whose main ingredients are:

\begin{enumerate}[topsep=0pt,itemsep=2pt]
  \item Using probability flux (rather than probability alone) to identify states
        that are dynamically important but may be rarely occupied (definitions in
        Section~3).
  \item Combining quantile-based pruning with a flux safeguard that prevents the
        removal of low-probability, high-flux bottleneck states (Section~3 and
        the flux-preserving pruning algorithm).
  \item Choosing time steps via a flux-based activity measure (implemented through
        a maximum-flux rule), which automatically shortens the step during stiff
        or transient phases and enlarges it in slow regimes (Section~3, with the
        rule motivated by the error analysis in Section~4).
  \item Developing a new error analysis based on a block-decomposed generator and
        boundary fluxes, leading to rigorous local and global $\ell^1$ error
        bounds via contraction of the stochastic semigroup (Proposition~4.1,
        Corollary~4.2, Theorem~4.4, and Corollaries~4.6--4.8 in Section~4).
  \item Demonstrating the algorithm on demanding benchmark problems, including a
        bottleneck system, a stochastic toggle switch, the oscillatory
        Oregonator, and the stiff Robertson (ROBER) autocatalytic system, with
        detailed reporting of algorithmic parameters, state-space sizes, and
        performance (Section~5, with the new bottleneck study in Section~5.1
        preceding the toggle switch and larger systems).
\end{enumerate}

Below we respond point by point to the reviewers’ comments.

%======================================================================
\section*{Response to Reviewer 1}
%======================================================================

\subsection*{Major Comments}

\noindent\textbf{Comment 1.1:} \textit{Lines 178--180: Can the authors elaborate why using a fixed probability
cutoff ``can be unreliable when the overall probability distribution changes
drastically over time''?}

\textbf{Response:} We have expanded both the conceptual explanation and the numerical
evidence for this point.

Conceptually, low probability does not imply negligible dynamical importance:
states with small $p(x,t)$ but large exit rate $w(x)$ can carry significant
probability flux and act as essential conduits between regions of state space.
This is now formalized in Section~3, where we define the outgoing probability flux
\[
  \Phi(x,t) = p(x,t)\,w(x)
\]
and the total system flux
\[
  \Phi_{\mathrm{total}}(J,t)
  = \sum_{x\in J} \Phi(x,t),
\]
and explicitly point out that such states can be connectivity-critical despite
being rarely occupied.

To make this concrete, the revised Introduction and Numerical Experiments section
use a bottleneck system inspired by your suggestion:
\[
  A \xrightarrow{k_1} B, \qquad B \xrightarrow{k_2} B + C,
\]
with $k_1 \ll k_2$ and initial condition $(N_A,N_B,N_C) = (1,0,0)$. In this system
the intermediate state $(N_A,N_B)=(0,1)$ has very small probability but large exit
flux, and it is the only pathway to increasing $C$.

Section~5.1 presents this bottleneck system in detail. The comparison there (Figure~4
and Tables~1--2 in the revised manuscript) shows that:
\begin{itemize}
  \item Under probability-only pruning (quantile step without the flux safeguard),
        states on this bottleneck path are removed as soon as their probability
        falls below the quantile threshold. The truncated generator becomes
        reducible and the dynamics are effectively trapped near $(1,0,0)$.
  \item Under the flux-preserving algorithm, the same states are protected
        because their $\Phi(x,t)$ constitutes a non-negligible fraction of the
        total system flux, even though $p(x,t)$ is tiny. As a result, the
        algorithm correctly captures the transition to large $C$ and the advancing
        probability front.
\end{itemize}

We hope this makes clear both why fixed probability cutoffs can be unreliable and
how the flux safeguard resolves this failure mode.

\medskip

\noindent\textbf{Comment 1.2:} \textit{Is it necessary to always renormalize the generator matrix and the
probability distribution? One powerful property of FSP comes from controlling
the loss of probability mass via $1-\mathbf{1}^\top\mathbf{p}(t)$. Also, if the
truncated generator is simply the principal submatrix of the full generator and
is never renormalized, then $1-\mathbf{1}^\top\mathbf{p}(t)$ is a global error
estimate. Can the authors elaborate why normalization is favorable or necessary
and why their more elaborate error control is better?}

\textbf{Response:} We appreciate this question and have clarified the distinction
between generator normalization and probability normalization, as well as the link
to the classical FSP bound.

\begin{itemize}
  \item \textbf{Generator.} In the revised method we do not renormalize the
        generator. Instead, Section~4 constructs a compressed generator
        $\widetilde{A}_{JJ}$ on the active set $J$ by folding transitions into and
        out of the complement $J'$ into a boundary term (see the definition of
        $\widetilde{A}_{JJ}$ in Section~4). As shown there,
        $\mathbf{1}^\top \widetilde{A}_{JJ} = 0$, so $\widetilde{A}_{JJ}$
        generates a contraction semigroup in $\ell^1$ (Proposition~4.1 and
        Corollary~4.2). This is conceptually similar to taking a principal
        submatrix, but it retains the effect of the truncated region through the
        boundary fluxes.

  \item \textbf{Probability.} The probability vector on $J$ is renormalized after
        each pruning event so that $\|\mathbf{p}_J(t_n)\|_1 = 1$ holds at the
        discrete times $t_n$. This keeps the numerical solution interpretable as a
        probability distribution on the current active set and improves numerical
        robustness when the remaining mass is small. This is particularly
        important because the active set is expanded and pruned many times over
        the course of a simulation, and maintaining this invariance simplifies
        both implementation and monitoring. In the analysis, the effect of
        renormalization appears explicitly in the local model error and is
        bounded by the pruned mass.

  \item \textbf{Relation to $1-\mathbf{1}^\top\mathbf{p}_J(t)$.} For a fixed FSP with
        generator equal to a principal submatrix of the full generator,
        $1-\mathbf{1}^\top\mathbf{p}_J(t)$ bounds the probability that has left
        $J$. Our block-partitioned analysis in Section~4 recovers and refines this
        view: we express the local model error over a time step $[t_n,t_{n+1}]$ in
        terms of boundary fluxes between $J$ and $J'$, and then show that the
        global error after $N$ steps is bounded by the sum of these local model
        errors plus the local time-stepping errors (Theorem~4.4 and
        Corollary~4.6).
\end{itemize}

Compared to a single scalar quantity $1-\mathbf{1}^\top\mathbf{p}_J(t)$, the
flux-based decomposition separates (i) model error due to truncation and pruning
(through boundary fluxes and pruned mass) from (ii) numerical error due to the
matrix exponential. This structure directly motivates our adaptive strategy,
which uses boundary fluxes to decide when to expand the state space and how to
adjust the time step. We have added a short remark in Section~4 to make this
connection explicit.

\medskip

\noindent\textbf{Comment 1.3:} \textit{The proof for Theorem 3.10 is incomplete. Specifically, I do not
understand why the errors could simply be summed over the $N$ time steps.}

\textbf{Response:} We agree that the earlier argument was incomplete. In the revised
manuscript, the old Theorem~3.10 has been removed and replaced by a new error
analysis in Section~4, culminating in an explicit global error bound.

The revised analysis proceeds as follows:

\begin{itemize}
  \item Proposition~4.1 and Corollary~4.2 show that the compressed generator
        $\widetilde{A}_{JJ}$ generates a contraction semigroup in $\ell^1$, i.e.,
        $\|\mathrm{e}^{\widetilde{A}_{JJ} t}\|_{1\to 1} = 1$ for $t\ge 0$.
  \item Theorem~4.4 bounds the local model error over one time step in terms of the
        integrated boundary fluxes across the truncation.
  \item A subsequent proposition bounds the local time-stepping error for the
        Krylov-based matrix exponential approximation, controlled by the ODE
        tolerance.
  \item Combining these pieces, Corollary~4.6 shows that the global $\ell^1$ error
        $e_n$ at time $t_n$ satisfies
        \[
          e_{n+1} \le e_n + \varepsilon_n + \tau_n,
        \]
        where $\varepsilon_n$ is the local model error and $\tau_n$ is the local
        time-stepping error. By induction this yields
        \[
          e_N \le \sum_{n=0}^{N-1} (\varepsilon_n + \tau_n).
        \]
\end{itemize}

The key point is that, because the compressed generator semigroup is contractive,
local errors do not get amplified as they propagate forward within a time step.
Errors from earlier steps therefore carry forward without growth and simply add to
the new local error. We now make this point explicit in the proof of the global
error bound in Section~4.

\medskip

\noindent\textbf{Comment 1.4:} \textit{Certain fringe models with ``bottleneck reactions'' might be
nontrivial with any FSP with time-stepping or sliding windows, including the
method in this paper.}

\textbf{Response:} We are very grateful for this comment; it directly motivated one of
the central case studies in the revised manuscript.

Section~5.1 now presents a detailed bottleneck example along the lines you
described. We consider a three-species system with a slow activation step and a
fast production step leading to an advancing front in $C$. The parameters are
chosen so that the intermediate states on the activation path have extremely low
probability but carry the only significant flux into regions with large $C$.

In this setting we compare:
\begin{itemize}
  \item a probability-only variant (quantile pruning without the flux safeguard), and
  \item the full flux-preserving algorithm (flux-based pruning combined with the
        flux-adaptive time-stepping rule of Section~3).
\end{itemize}

The results in Section~5.1 (Figure~4 and Tables~1--2) show that the
probability-only variant quickly removes bottleneck states from the truncated set,
after which the numerical solution remains localized near the initial state and
fails to reproduce the growth of $C$. In contrast, the flux-based variant
retains bottleneck states whenever their exit flux is a non-negligible fraction
of the total system flux, and it reproduces the expected increase in
$\langle C\rangle$ over several orders of magnitude in time. This example
illustrates precisely the pathology you raised and shows that the flux safeguard
is effective in protecting such bottleneck channels.

\medskip

\noindent\textbf{Comment 1.5:} \textit{Can the authors compare their method to previous FSP methods that
also incorporate adaptive state spaces, either in discussion or numerical tests?
Some are already cited in the paper. Another paper that comes to mind is Dinh \&
Sidje (Physical Biology, 2017).}

\textbf{Response:} We have expanded the discussion of related work in the Introduction
(Related Work subsection) and in Section~5 to position the revised method more
clearly among existing approaches.

In particular, we now discuss:
\begin{itemize}
  \item Aggregation and lumping methods such as those of Bobbio \& Trivedi and
        Zhang, Watson, \& Cao, which group micro-states into macro-states using
        stationary or Monte Carlo information.
  \item Sliding-window and hybrid FSP--SSA schemes such as those by Wolf
        et al.\ and MacNamara et al., which track moving regions of support or
        couple deterministic and stochastic updates.
  \item Krylov-FSP-SSA and tensor-train methods of Dinh \& Sidje and
        collaborators, which guide state-space expansion using SSA trajectories or
        compressed representations.
\end{itemize}

Our flux-based adaptive FSP differs in two main ways:

\begin{enumerate}
  \item In the analysis, we control truncation error using flux across the
        active--inactive boundary, which leads directly to the local and global
        error bounds in Section~4. In the algorithm, we approximate this boundary
        activity using statewise exit flux as the main control quantity rather
        than relying on stationary information.
  \item We decouple discovery of new states (via reachability expansion of the
        reaction network) from the pruning decision, using quantiles only as a
        pre-selection mechanism and flux as the final safeguard. This explicitly
        addresses the bottleneck scenario highlighted by Reviewer~1.
\end{enumerate}

Implementing and tuning several competing adaptive FSP codes for a full numerical
benchmark on all our examples would require substantial additional effort and is
beyond the scope of the current revision. Instead, we focus on clearly articulating
the conceptual differences and providing detailed algorithmic descriptions and
performance data that we hope will make future comparisons easier.

\medskip

\noindent\textbf{Comment 1.6:} \textit{Although well studied, Michaelis--Menten, Lotka--Volterra, and toggle
switch models are relatively simple to solve. It would be beneficial to showcase
your algorithm with at least one example where the state space is known to be
large. Can the authors implement such a model?}

\textbf{Response:} We agree that more demanding examples are essential. The revised
manuscript now contains several such systems and explicitly orders them so that
the bottleneck example and its error analysis appear first in Section~5.1, followed by
the toggle switch and higher-dimensional models.

\begin{itemize}
  \item Section~5.1 presents the bottleneck reaction network and, in the
        ``Experimental Error Analysis'' subsubsection, compares the flux-based
        adaptive FSP to a high-accuracy fixed-FSP reference on the bounded state
        space $\{0,\dots,2\}\times\{0,\dots,2\}\times\{0,\dots,10000\}$. Figure~5
        shows that the relative error in $\langle C\rangle$ remains below $0.1\%$
        and decreases to about $0.015\%$ as the distribution spreads, with an
        approximately constant instantaneous error rate in time, in agreement with
        the additive error structure proved in Section~4. Table~3 reports mean
        values and errors for all species at $t=10^5$ s; the adaptive and
        reference solutions are in close agreement, with the adaptive state space
        about $4.5$ times smaller than the fixed one.
  \item Section~5.2 revisits the stochastic toggle switch, now primarily as a
        qualitative illustration of bistability and basin switching under
        flux-preserving truncation (Figure~6).
  \item Section~5.3 introduces the Oregonator model of chemical oscillations.
        Based on observed mean trajectory ranges, a conservative rectangular
        bounding box would contain on the order of
        $4{,}000\times 7{,}000\times 11{,}000 \approx 3\times 10^{11}$ states,
        whereas the adaptive FSP keeps between a few dozen and a few thousand
        states active at any time, achieving a compression ratio of roughly
        $10^8{:}1$ while faithfully reproducing the oscillatory dynamics
        (Figure~7).
  \item Section~5.4 adds the Robertson (ROBER) autocatalytic system, which
        is extremely stiff, with reaction rates spanning nine orders of magnitude.
        Using the conservation law $N_A+N_B+N_C=10{,}000$ and the observed ranges
        for $N_A$ and $N_B$, a conservative rectangular bounding box would contain
        on the order of $10^5$ states, whereas the adaptive state space in our
        simulations remains in the range $50$--$2{,}000$ states. Thus the method
        achieves a compression of about two orders of magnitude while capturing
        the three-phase dynamics (slow initiation, explosive autocatalysis, and
        long equilibration) and controlling local boundary flux and time-stepping
        error (Figure~8).
\end{itemize}


%======================================================================
\section*{Response to Reviewer 2}
%======================================================================

\subsection*{Main Comments}

\noindent\textbf{Comment 2.1:} \textit{The first experiment, depicted in Figure 2, requires a much deeper
discussion. The error seems to increase with time, contradicting the analysis.
Also, 1000 SSA trajectories would not be enough to eliminate statistical error.}

\textbf{Response:} We agree that the original SSA-based experiment did not clearly
separate truncation error from Monte Carlo sampling error, and that the discussion
of error growth was incomplete.

In the revised manuscript we have replaced this experiment with a deterministic
error study that directly tests the theoretical results:

\begin{itemize}
  \item Section~5.1.1 (Experimental Error Analysis) compares the flux-based
        adaptive FSP to a high-accuracy fixed-FSP reference on the bounded state
        space $\{0,\dots,2\}\times\{0,\dots,2\}\times\{0,\dots,10000\}$.
  \item Figure~5 shows that the relative error in $\langle C\rangle$ remains below
        $0.1\%$ and decreases to about $0.015\%$ as the distribution spreads, with
        an approximately constant instantaneous error rate in time, in agreement
        with the additive error structure in Section~4.
  \item Table~3 reports mean values and errors for all species at $t=10^5$ s, with
        the adaptive and reference solutions in close agreement and the adaptive
        state space about $4.5$ times smaller than the fixed one.
\end{itemize}

We believe this new experiment more directly and transparently validates the
revised error analysis and avoids conflating truncation error with statistical
sampling error.

\medskip

\noindent\textbf{Comment 2.2:} \textit{For most (if not all) experiments, the pruning factor $\alpha$, the
effective pruned fraction $m$, and the time step are not reported. These should
be specified in the paper rather than needing to be inferred from the code.}

\textbf{Response:} We have revised the numerical sections to state the key algorithmic
parameters explicitly:

\begin{itemize}
  \item Section~3, which introduces the flux-preserving pruning algorithm, now
        specifies the quantile tolerance $\alpha$, the flux tolerance
        $\varepsilon_{\rm flux}$, and the bounds $\Delta t_{\min}$ and
        $\Delta t_{\max}$ used in the bottleneck experiments, as well as the
        functional form of the flux-based time-step rule.
  \item Section~5.1.1 provides the time grid and state-space bounds for the
        fixed-FSP reference solution, and the adaptive time-step settings for the
        comparison run in the bottleneck error study.
  \item The later subsections of Section~5 (toggle switch, Oregonator, and ROBER)
        state the choices of $\alpha$, flux tolerance, and time-step range used in
        each example. The figures also show the resulting adaptive time steps and
        state-space sizes as functions of time.
\end{itemize}

In the revised method, $\alpha$ primarily determines how much probability mass is
eligible for pruning per step, while the flux safeguard further constrains which
of those states can actually be removed. Rather than reporting an effective $m$ at
every time step, we focus on how $\alpha$, the flux threshold, and the boundary
fluxes enter the local and global error bounds in Section~4.

\medskip

\noindent\textbf{Comment 2.3:} \textit{I would like to see more explicitly the relation between the pruning
factor $\alpha$ and the effective pruned fraction $m$. This relationship will
affect both efficiency and accuracy.}

\textbf{Response:} We appreciate this request and have clarified the role of $\alpha$
in the revised algorithm.

In the original quantile-only scheme, $\alpha$ directly controlled the fraction of
probability mass removed at each pruning step, so $m \approx \alpha$. The new
method introduces a flux-based safeguard:

\begin{itemize}
  \item The quantile parameter $\alpha$ selects a candidate set of low-probability
        states whose cumulative probability does not exceed $\alpha$.
  \item The flux safeguard then scans this candidate set and retains any state
        whose outgoing flux exceeds a prescribed fraction of the total system
        flux. These protected states are not pruned, regardless of their
        probability.
\end{itemize}

As a consequence, the effective pruned mass $m$ is at most $\alpha$, and can
be considerably smaller when bottleneck states are present. This behavior is
illustrated in the bottleneck experiment (Section~5.1): even with a relatively
large $\alpha$, the flux safeguard prevents removal of the intermediate bottleneck
states and thereby limits $m$ while maintaining connectivity.

The local error analysis in Section~4 makes explicit how the local model error
depends on the pruned mass, the boundary fluxes, and the time step, thereby
connecting $\alpha$ and $m$ through their impact on the flux-based error control.
A full parametric study of $(\alpha,m)$ over many systems would significantly
lengthen the paper; we therefore leave a more systematic investigation to future
work.

\medskip

\noindent\textbf{Comment 2.4:} \textit{The authors report runtimes for their new algorithm, but these are
only average values on a single hardware platform. At least the hardware used
should be mentioned, and runtimes should be compared to those of existing
algorithms; probably higher-dimensional examples need to be studied as well.}

\textbf{Response:} We agree that performance information and reproducibility are
important. In the revised manuscript we have:

\begin{itemize}
  \item Specified the implementation details and hardware in the Numerical
        Experiments section: all experiments were run in Julia using Expokit.jl
        for Krylov-based matrix exponential computations, Catalyst.jl for
        reaction-network modeling, and sparse matrices from the
        \texttt{SparseArrays} standard library. Simulations were performed on a
        MacBook Pro with an M1 Pro chip and 16\,GB of RAM.
  \item Added a dedicated subsection on performance evaluation (Section~5.5), which
        analyzes the performance of forward-enumeration matrix construction
        compared to a standard backward-expansion strategy. At the kernel level,
        benchmark plots show that, for the Oregonator model, forward enumeration
        reduces the median time for generator construction from about
        $8.91$\,ms to $0.62$\,ms, a speedup of roughly $14.4\times$. In terms of
        total simulation cost, Figure~10 and Table~8 report that forward
        enumeration yields speedups of about $3.9\times$ (Oregonator),
        $3.5\times$ (ROBER), and $9.8\times$ (toggle switch) relative to a
        reference implementation using backward expansion.
\end{itemize}

A direct head-to-head comparison with existing adaptive FSP implementations
(e.g., Optimal FSP, Krylov-FSP-SSA, aggregation-based methods) on the same
hardware would require reimplementing or interfacing with multiple external codes
and carefully tuning each method for fairness. We view such a study as an
interesting direction for future work and, for the present revision, focus on
documenting the internal performance characteristics of the proposed algorithm.

\medskip

\noindent\textbf{Comment 2.5:} \textit{Include a test case where excessive pruning would be an issue but
your new method avoids that problem.}

\textbf{Response:} The bottleneck example in Section~5.1 was designed precisely to
address this request (and the related concerns raised by Reviewer~1). This example
became the starting point for shifting the scope of the paper toward stiff,
multiscale systems where rare but dynamically essential transitions must be
preserved.

As described in our response to Comment~1.4, the bottleneck model contains
low-probability states that form the only pathway to large $C$. Any scheme that
prunes solely on probability is forced either to choose a very small $\alpha$
(sacrificing efficiency) or eventually to remove these states and thus
qualitatively change the dynamics. The flux-based adaptive FSP method, by contrast,
monitors the outgoing flux from candidate states and protects those whose flux is
non-negligible. The plots and tables in Section~5.1 show that the flux-based
variant retains the bottleneck channel and recovers the correct long-time
behavior, whereas the probability-only variant does not.

%======================================================================
\section*{Summary}
%======================================================================

To summarize, the revised manuscript differs substantially from the original
submission. The main advances are:

\begin{enumerate}
  \item A flux-preserving adaptive FSP algorithm that augments quantile-based
        pruning with a flux safeguard and a flux-guided time-step controller,
        directly addressing the bottleneck failure modes pointed out by
        Reviewer~1.
  \item A new error analysis that expresses local model error in terms of
        boundary fluxes and exploits semigroup contraction to obtain additive
        global error bounds (Proposition~4.1, Corollary~4.2, Theorem~4.4, and
        Corollaries~4.6--4.8 in Section~4). In the Discussion we also outline a
        natural next step of relating our flux-based error control to spectral
        quantities (e.g., spectral gaps) via ergodicity coefficients and
        Cheeger-type inequalities.
  \item A set of challenging numerical experiments, including a detailed
        bottleneck case study and its error analysis in Section~5.1, a stochastic
        toggle switch, the oscillatory Oregonator, and the stiff Robertson system,
        each with explicit reporting of parameters, state-space sizes, and
        observed behavior (Figures~4--8, Tables~1--3).
  \item A performance study of forward-enumeration matrix construction that
        demonstrates substantial speedups over a standard backward-expansion
        approach (Figure~10 and Table~8 in Section~5.5).
\end{enumerate}

We hope that these changes fully address the reviewers’ concerns and that the
revised manuscript is now suitable for publication in \emph{Multiscale Modeling
and Simulation}. We are very grateful to both reviewers and to the Associate
Editor for their detailed and thoughtful feedback, which has been instrumental in
improving this work.

\vspace{1em}
\noindent Sincerely,\\[0.5em]
Aditya Dendukuri, Shivkumar Chandrasekaran, and Linda Petzold

\end{document}